%-------------------------
% Cover letter - Miovision, Senior Software Engineer (Platform Engineering)
% Author : Aymen Daoudi
% Reuses Parameters.sty / Commands.sty from the resume so both documents match.
%------------------------

\documentclass[letterpaper,10pt]{article}

% Parameters.sty, Commands.sty and Images/ live in the repo root, one level up.
% \input@path lets \usepackage resolve them from inside this folder.
\makeatletter
\def\input@path{{../}}
\makeatother

\usepackage{Parameters}
\usepackage{Commands}

% graphicx is loaded by Parameters.sty, so \graphicspath must come after it.
% This resolves the header icons under Images/.
\graphicspath{{../}}

% Improves inter-word spacing in justified text; keeps the rivers down.
\usepackage{microtype}

% Letter geometry: the resume's 1cm margins are too tight for running prose.
\geometry{left=2.6cm, right=2.6cm, top=1.6cm, bottom=1.6cm}

% Block paragraph style, no indentation.
\setlength{\parindent}{0pt}
\setlength{\parskip}{7pt}
\setlength{\footskip}{12pt}
\linespread{1.05}

\begin{document}

% \roboto is a declaration, not a wrapper. Commands.sty issues it per-macro, so
% the resume renders in Roboto; running prose here needs it set once, document
% wide, or it falls back to the Latin Modern serif default.
\roboto

% Parameters.sty issues a global \raggedright for the resume layout. Undo it so
% the letter body is justified. \parindent stays at 0pt for block paragraphs.
\rightskip=0pt
\leftskip=0pt
\parfillskip=0pt plus 1fil

%----------HEADING (mirrors Parts/Heading.tex)-----------------

\begin{tabular*}{\textwidth}{l@{\extracolsep{\fill}}l}
  \multirow{2}{*}{\Name{\huge A\LARGE YMEN \huge D\LARGE AOUDI}}
    & \Addresse{\quad Montréal, Québec, Canada} \\
    & \EmailAddresse{\quad \href{mailto:iaymendaoudi@gmail.com}{iaymendaoudi@gmail.com}} \\
  \RoleDescription{Software Development Engineer}
    & \PhoneNumbers{\quad +1 438 927 9580}{} \\
    & \faLinkedin \enspace \enspace {\url{https://www.linkedin.com/in/aymendaoudi}}
\end{tabular*}

\vspace{6pt}
{\color{black}\rule{\textwidth}{0.4pt}}
\vspace{10pt}

%----------LETTER-----------------

{\footnotesize August 22, 2026}

\vspace{2pt}

Dear Miovision Hiring Team,

I am applying for the Senior Software Engineer position on the Platform
Engineering team. The line that decided it for me is the one about designing
machine-readable API contracts that AI agents can discover and invoke safely.
That is the problem I work on now from the other side of it, and not many backend
roles are treating it as a first-class design concern yet.

I am a Senior Software Engineer at Numa, where an LLM-driven phone agent
completes real transactions by calling into systems we do not own. My work is the
layer that makes that safe: the integrations with dealer management systems,
including the Reynolds \& Reynolds one I led, and the contracts governing what the
agent may invoke on its own, what it has to hand back to a person, and how
failures surface. I have also had generative models running in a live call flow,
where latency and failure handling matter more than model capability, and my own
development runs on AI-assisted tooling in Cursor and Claude Code.

Before Numa I spent three years at Divisions Maintenance Group, redesigning
services around domain-driven design and Onion architecture, building the
internal tooling that scaffolded new APIs consistently, and moving teams onto a
release flow that cleared a bottleneck they had lived with for a long time. That
mix of API design, developer tooling and cross-team agreement is the closest
analogue I have to platform work. Earlier roles covered maritime compliance
software at Lloyd's Register OneOcean, electric vehicle charging at AddEnergie,
and financial policy and insurance at Stratec Global, on microservices over
Docker, Kubernetes, AWS and Azure.

What draws me to Miovision is that this gateway is being built rather than
retrofitted, which is when domain modelling and governance decisions actually
hold. That is where I would contribute: API and domain design, turning cross-team
disagreement into contracts people keep to, and the agent-facing side of the
platform. I should also be straight about one gap. I have not written Kotlin
professionally. My backend work is C\#, Python and TypeScript, and my JVM
experience is Java, which I taught at UQAM and used for my master's research on
static analysis of Android bytecode. I have changed stacks several times and
would not be starting from nothing, but that is better weighed now than
discovered later.

Thank you for considering my application. I would be glad to discuss the role and
where I would fit.

\vspace{6pt}

Respectfully,

\vspace{4pt}

Aymen Daoudi

\end{document}
